% ===========================================================================
% Part IV - Ensemble Learning (Experiments 8-12)
% ===========================================================================

\part{Ensemble Learning}

% ---------------------------------------------------------------------------
\chapter{Random Forest Regression (RFR)}\label{exp:08}

\runninghead{Experiment 8: Random Forest Regression (RFR)}

\quickfacts
  {Dataset}{California Housing, 5,000-row random sample, 8 features; 80/20 split}
  {Key parameters}{150 trees, max\_depth = 12, out-of-bag scoring enabled}
  {Headline result}{Test $R^2$ 0.7415 (OOB 0.7513); RMSE 0.5908; MAE 0.3999}

\section{Objective}
Predict a continuous target using an ensemble of randomized decision trees; evaluate with
MAE, RMSE and $R^2$ and inspect out-of-bag validation.

\section{Suggested Dataset}
California Housing (scikit-learn), 5,000-row random sample, target = median house value.

\section{Required Libraries}
scikit-learn, numpy, pandas.

\section{Brief Theory}
Random Forest Regression trains $B$ decision trees on bootstrap samples of the training data;
at each split only a random subset of features is considered, which decorrelates the trees.
The ensemble prediction is the average of the tree outputs,
\[
\hat{f}_{\mathrm{RF}}(x) = \frac{1}{B}\sum_{b=1}^{B} T_b(x) .
\]
Averaging approximately uncorrelated trees reduces variance: for pairwise correlation $\rho$,
$\mathrm{Var}(\hat{f}_{\mathrm{RF}}) = \rho\sigma^2 + \frac{1-\rho}{B}\sigma^2$, so adding
trees can only help. Out-of-bag (OOB) samples --- the $\approx 37\%$ not drawn in a bootstrap
--- provide a free internal validation estimate.

\section{Procedure}
\begin{enumerate}
  \item Load and sample the dataset; split into training and test sets.
  \item Train \texttt{RandomForestRegressor} with 150 trees, \texttt{max\_depth} = 12 and
        OOB scoring enabled.
  \item Predict the test set and compute $R^2$, RMSE and MAE.
  \item Report the OOB $R^2$ and compare it with the test $R^2$ to check overfitting.
  \item Inspect feature importances.
\end{enumerate}

\section{Reference Implementation}
% ===========================================================================
% exp08_rfr.tex - complete code listings for Experiment 8
% Source: notebooks/04_ensemble_learning.ipynb (executed)
% Repository: https://github.com/Atik203/ML-Algorithms
% ===========================================================================

\begin{codebox}[title={Core libraries}]
# ---- Core libraries ----
import numpy as np
import pandas as pd
import matplotlib.pyplot as plt
import seaborn as sns

from sklearn.datasets import fetch_california_housing
from sklearn.model_selection import train_test_split
from sklearn.ensemble import RandomForestRegressor, RandomForestClassifier, AdaBoostRegressor, AdaBoostClassifier
from sklearn.metrics import r2_score, root_mean_squared_error, mean_absolute_error, accuracy_score, f1_score, roc_auc_score

import xgboost as xgb
import catboost as cb

# Styling setup
plt.style.use('seaborn-v0_8-whitegrid' if 'seaborn-v0_8-whitegrid' in plt.style.available else 'default')
plt.rcParams['font.sans-serif'] = 'DejaVu Sans'
plt.rcParams['figure.figsize'] = (10, 6)
plt.rcParams['figure.dpi'] = 110
np.random.seed(42)
\end{codebox}

\begin{codebox}[title={Load the California Housing dataset (real-world regression benchmark)}]
# ---- Load the California Housing dataset (real-world regression benchmark) ----
housing = fetch_california_housing(as_frame=True)
# Sample 5,000 rows so every ensemble trains quickly while remaining representative
df_reg = housing.frame.sample(n=5000, random_state=42)

X_reg = df_reg[housing.feature_names]   # 8 socio-economic / geographic features
y_reg = df_reg['MedHouseVal']           # target: median house value

# ---- Train/test split (80% / 20%) ----
X_train_reg, X_test_reg, y_train_reg, y_test_reg = train_test_split(
    X_reg, y_reg, test_size=0.20, random_state=42
)

print(f"Regression Training Samples: {X_train_reg.shape[0]}, Test Samples: {X_test_reg.shape[0]}")
df_reg.head()
\end{codebox}

\begin{codebox}[title={Random Forest Regression (bagging of decision trees)}]
# ---- Random Forest Regression (bagging of decision trees) ----
# n_estimators = number of trees, oob_score = estimate accuracy from out-of-bag samples
rf_reg = RandomForestRegressor(n_estimators=150, max_depth=12, oob_score=True, random_state=42, n_jobs=-1)
rf_reg.fit(X_train_reg, y_train_reg)

# ---- Predict and evaluate ----
y_pred_rf = rf_reg.predict(X_test_reg)
r2_rf = r2_score(y_test_reg, y_pred_rf)                          # explained variance (1.0 = perfect)
rmse_rf = root_mean_squared_error(y_test_reg, y_pred_rf)         # error in target units
mae_rf = mean_absolute_error(y_test_reg, y_pred_rf)              # robust average error

print(f"Random Forest Regressor Results:")
print(f" - Out-of-Bag (OOB) R2 Score: {rf_reg.oob_score_:.4f}")
print(f" - Test R2 Score:              {r2_rf:.4f}")
print(f" - Test RMSE:                  {rmse_rf:.4f}")
print(f" - Test MAE:                   {mae_rf:.4f}")
\end{codebox}


\section{Evaluation / Expected Output}
\begin{table}[htbp]
  \centering
  \caption{Random Forest Regression on the California Housing sample.}
  \begin{tabular}{|l|r|}
    \hline
\rowcolor{mlblue!8}    \textbf{Metric} & \textbf{Value} \\
    \hline
    OOB $R^2$ & 0.7513 \\
\hline
    Test $R^2$ & 0.7415 \\
\hline
    Test RMSE & 0.5908 (in \$100k units $\approx$ \$59k) \\
\hline
    Test MAE & 0.3999 (in \$100k units $\approx$ \$40k) \\
    \hline
  \end{tabular}
\end{table}

The OOB $R^2$ (0.7513) is very close to the test $R^2$ (0.7415), so the model is not
overfitting and the OOB estimate is a reliable free validation. Feature importances are
dominated by median income and location (latitude/longitude), consistent with housing
economics.

\section{Viva / Discussion Questions}
\begin{enumerate}
  \item \textbf{Why does a forest reduce variance?}\par\nobreak Averaging $B$ approximately uncorrelated trees
  gives $\rho\sigma^2 + (1-\rho)\sigma^2/B$; as $B$ grows the second term vanishes, so
  predictions are more stable than a single tree's.
  \item \textbf{What is bagging?}\par\nobreak Bootstrap aggregating: each model is trained on a bootstrap sample
  of the data and the predictions are combined (average for regression, vote for
  classification).
  \item \textbf{Why can feature importance be misleading?}\par\nobreak Impurity-based importance is biased toward
  high-cardinality/continuous features, is computed on training data (it implies no causality),
  and correlated features share or dilute importance arbitrarily.
\item \textbf{Why does adding more trees not make a random forest overfit?}\par\nobreak Averaging many decorrelated trees reduces variance without increasing bias; beyond a certain number the error stabilizes, so more trees cost computation but not generalization.
\end{enumerate}


% ---------------------------------------------------------------------------
\chapter{Random Forest Classification (RFC)}\label{exp:09}

\runninghead{Experiment 9: Random Forest Classification (RFC)}

\quickfacts
  {Dataset}{Heart Disease (\texttt{data/heart\_disease.csv}, 303 rows, 13 features); stratified 75/25 split}
  {Key parameters}{120 trees, max\_depth = 8; one-hot encoded categorical features}
  {Headline result}{Accuracy 0.7763; F1 0.8046; ROC-AUC 0.8704}

\section{Objective}
Classify observations using an ensemble of randomized decision trees, with a stratified
split and classification metrics.

\section{Suggested Dataset}
Heart Disease (\texttt{data/heart\_disease.csv}), 303 rows, 13 features plus the binary target.

\section{Required Libraries}
scikit-learn, numpy, pandas.

\section{Brief Theory}
Random Forest Classification aggregates tree votes. Each tree votes for one class and the
ensemble predicts the most frequent class,
\[
\hat{C}_{\mathrm{RF}}(x) = \arg\max_k \sum_{b=1}^{B} \mathbb{I}\bigl(T_b(x) = k\bigr) .
\]
Bootstrap sampling plus feature subsampling decorrelates the trees, reducing variance relative
to a single decision tree. The fraction of trees voting for the positive class also provides a
class probability, which is used for ROC-AUC.

\section{Procedure}
\begin{enumerate}
  \item Load the heart data; one-hot encode the categorical features.
  \item Use a stratified 75/25 split to preserve the disease ratio.
  \item Train \texttt{RandomForestClassifier} with 120 trees and \texttt{max\_depth} = 8.
  \item Predict labels and probabilities on the test set.
  \item Compute accuracy, F1 and ROC-AUC.
\end{enumerate}

\section{Reference Implementation}
% ===========================================================================
% exp09_rfc.tex - complete code listings for Experiment 9
% Source: notebooks/04_ensemble_learning.ipynb (executed)
% ===========================================================================

\begin{codebox}[title={Core libraries}]
# ---- Core libraries ----
import numpy as np
import pandas as pd
import matplotlib.pyplot as plt
import seaborn as sns

from sklearn.datasets import fetch_california_housing
from sklearn.model_selection import train_test_split
from sklearn.ensemble import RandomForestRegressor, RandomForestClassifier, AdaBoostRegressor, AdaBoostClassifier
from sklearn.metrics import r2_score, root_mean_squared_error, mean_absolute_error, accuracy_score, f1_score, roc_auc_score

import xgboost as xgb
import catboost as cb

# Styling setup
plt.style.use('seaborn-v0_8-whitegrid' if 'seaborn-v0_8-whitegrid' in plt.style.available else 'default')
plt.rcParams['font.sans-serif'] = 'DejaVu Sans'
plt.rcParams['figure.figsize'] = (10, 6)
plt.rcParams['figure.dpi'] = 110
np.random.seed(42)
\end{codebox}

\begin{codebox}[title={Load the Heart Disease dataset (real-world classification benchmark)}]
# ---- Load the Heart Disease dataset (real-world classification benchmark) ----
df_clf = pd.read_csv('../data/heart_disease.csv')
print(f"Heart Disease dataset: {df_clf.shape[0]} rows x {df_clf.shape[1]} columns")

# Separate target and features (handle either 'target' or 'Target' as the column name)
target_col = 'target' if 'target' in df_clf.columns else 'Target'
y_clf = df_clf[target_col].values
X_clf_raw = df_clf.drop(columns=[target_col])

# One-hot encode categorical variables for RFC / XGBoost / AdaBoost
X_clf_encoded = pd.get_dummies(X_clf_raw, drop_first=True)

# Stratified split keeps the disease/no-disease ratio identical in train and test
X_train_clf, X_test_clf, y_train_clf, y_test_clf = train_test_split(
    X_clf_encoded, y_clf, test_size=0.25, random_state=42, stratify=y_clf
)

print(f"Classification Training Set: {X_train_clf.shape}, Test Set: {X_test_clf.shape}")

# ---- Random Forest Classification (majority vote of trees) ----
rfc = RandomForestClassifier(n_estimators=120, max_depth=8, random_state=42, n_jobs=-1)
rfc.fit(X_train_clf, y_train_clf)

# Predictions and probability of the positive (disease) class
y_pred_rfc = rfc.predict(X_test_clf)
y_prob_rfc = rfc.predict_proba(X_test_clf)[:, 1]

# ---- Evaluate: accuracy, F1 and ROC-AUC ----
acc_rfc = accuracy_score(y_test_clf, y_pred_rfc)
f1_rfc = f1_score(y_test_clf, y_pred_rfc)
auc_rfc = roc_auc_score(y_test_clf, y_prob_rfc)

print(f"Random Forest Classifier: Accuracy = {acc_rfc:.4f}, F1 = {f1_rfc:.4f}, ROC-AUC = {auc_rfc:.4f}")
\end{codebox}


\section{Evaluation / Expected Output}
\begin{table}[htbp]
  \centering
  \caption{Random Forest Classification on Heart Disease (76 test samples).}
  \begin{tabular}{|l|r|}
    \hline
\rowcolor{mlblue!8}    \textbf{Metric} & \textbf{Value} \\
    \hline
    Accuracy & 0.7763 \\
\hline
    F1-score & 0.8046 \\
\hline
    ROC-AUC & 0.8704 \\
    \hline
  \end{tabular}
\end{table}

ROC-AUC 0.87 indicates good ranking quality on a small dataset; accuracy 0.776 is typical for
heart-disease prediction with these features and reflects the problem's noise rather than a
modelling failure. F1 and AUC are safer summaries than accuracy alone because the classes are
not perfectly balanced.

\section{Viva / Discussion Questions}
\begin{enumerate}
  \item \textbf{RFR vs RFC?}\par\nobreak Regression averages the trees' numeric outputs; classification takes a
  majority vote (and averages class probabilities for scoring).
  \item \textbf{What is majority voting?}\par\nobreak Each tree votes for one class and the most frequent class
  becomes the ensemble prediction; soft voting averages the probabilities instead.
  \item \textbf{How can class imbalance affect accuracy?}\par\nobreak With a rare positive class, a model can
  predict the majority class everywhere and still show high accuracy; precision, recall, F1 and
  AUC expose the failure.
\item \textbf{Why is ROC-AUC threshold-independent?}\par\nobreak It summarizes performance over all classification thresholds by comparing true- and false-positive rates, so it does not depend on a single cut-off.
\end{enumerate}


% ---------------------------------------------------------------------------
\chapter{XGBoost}\label{exp:10}

\runninghead{Experiment 10: XGBoost}

\quickfacts
  {Dataset}{Heart Disease (same stratified 75/25 split as RFC)}
  {Key parameters}{120 trees, learning rate 0.08, max\_depth = 4, subsample / colsample = 0.8}
  {Headline result}{Accuracy 0.7763; F1 0.8132; ROC-AUC 0.8509}

\section{Objective}
Train a gradient-boosted tree classifier and analyze the effect of boosting
hyperparameters.

\section{Suggested Dataset}
Heart Disease (\texttt{data/heart\_disease.csv}), the same stratified split used for the
other ensembles.

\section{Required Libraries}
xgboost, scikit-learn, numpy, pandas.

\section{Brief Theory}
XGBoost builds trees sequentially, each fitting a second-order Taylor expansion of the loss:
\[
\mathcal{L}^{(t)} \approx \sum_{i=1}^{n}\Bigl[g_i f_t(x_i) + \tfrac{1}{2} h_i f_t^2(x_i)\Bigr] + \Omega(f_t),
\qquad
\Omega(f) = \gamma T + \tfrac{1}{2}\lambda \sum_{j=1}^{T} w_j^2 ,
\]
where $g_i$ and $h_i$ are the first and second derivatives of the loss with respect to the
current prediction. The regularization term $\Omega$ penalizes tree complexity and leaf
weights, while row/column subsampling adds stochasticity; both control overfitting.

\section{Procedure}
\begin{enumerate}
  \item Prepare the one-hot encoded heart data and the same stratified split.
  \item Fit \texttt{XGBClassifier} with 120 trees, learning rate 0.08, depth 4, subsample and
        column sample 0.8, and the logistic-loss objective.
  \item Predict labels and probabilities on the test set.
  \item Report accuracy, F1 and ROC-AUC.
  \item Discuss the complexity/generalization trade-off.
\end{enumerate}

\section{Reference Implementation}
% ===========================================================================
% exp10_xgboost.tex - complete code listings for Experiment 10
% Source: notebooks/04_ensemble_learning.ipynb (executed)
% ===========================================================================

\begin{codebox}[title={Core libraries}]
# ---- Core libraries ----
import numpy as np
import pandas as pd
import matplotlib.pyplot as plt
import seaborn as sns

from sklearn.datasets import fetch_california_housing
from sklearn.model_selection import train_test_split
from sklearn.ensemble import RandomForestRegressor, RandomForestClassifier, AdaBoostRegressor, AdaBoostClassifier
from sklearn.metrics import r2_score, root_mean_squared_error, mean_absolute_error, accuracy_score, f1_score, roc_auc_score

import xgboost as xgb
import catboost as cb

# Styling setup
plt.style.use('seaborn-v0_8-whitegrid' if 'seaborn-v0_8-whitegrid' in plt.style.available else 'default')
plt.rcParams['font.sans-serif'] = 'DejaVu Sans'
plt.rcParams['figure.figsize'] = (10, 6)
plt.rcParams['figure.dpi'] = 110
np.random.seed(42)
\end{codebox}

\begin{codebox}[title={Load the Heart Disease dataset (real-world classification benchmark)}]
# ---- Load the Heart Disease dataset (real-world classification benchmark) ----
df_clf = pd.read_csv('../data/heart_disease.csv')
print(f"Heart Disease dataset: {df_clf.shape[0]} rows x {df_clf.shape[1]} columns")

# Separate target and features (handle either 'target' or 'Target' as the column name)
target_col = 'target' if 'target' in df_clf.columns else 'Target'
y_clf = df_clf[target_col].values
X_clf_raw = df_clf.drop(columns=[target_col])

# One-hot encode categorical variables for RFC / XGBoost / AdaBoost
X_clf_encoded = pd.get_dummies(X_clf_raw, drop_first=True)

# Stratified split keeps the disease/no-disease ratio identical in train and test
X_train_clf, X_test_clf, y_train_clf, y_test_clf = train_test_split(
    X_clf_encoded, y_clf, test_size=0.25, random_state=42, stratify=y_clf
)

print(f"Classification Training Set: {X_train_clf.shape}, Test Set: {X_test_clf.shape}")

# ---- Random Forest Classification (majority vote of trees) ----
rfc = RandomForestClassifier(n_estimators=120, max_depth=8, random_state=42, n_jobs=-1)
rfc.fit(X_train_clf, y_train_clf)

# Predictions and probability of the positive (disease) class
y_pred_rfc = rfc.predict(X_test_clf)
y_prob_rfc = rfc.predict_proba(X_test_clf)[:, 1]

# ---- Evaluate: accuracy, F1 and ROC-AUC ----
acc_rfc = accuracy_score(y_test_clf, y_pred_rfc)
f1_rfc = f1_score(y_test_clf, y_pred_rfc)
auc_rfc = roc_auc_score(y_test_clf, y_prob_rfc)

print(f"Random Forest Classifier: Accuracy = {acc_rfc:.4f}, F1 = {f1_rfc:.4f}, ROC-AUC = {auc_rfc:.4f}")
\end{codebox}

\begin{codebox}[title={XGBoost classification on the same stratified split}]
# ---- XGBoost classification on the same stratified split ----
xgb_clf = xgb.XGBClassifier(
    n_estimators=120, learning_rate=0.08, max_depth=4,
    subsample=0.8, colsample_bytree=0.8,
    random_state=42, eval_metric='logloss'
)
xgb_clf.fit(X_train_clf, y_train_clf)

y_pred_xgb_c = xgb_clf.predict(X_test_clf)
y_prob_xgb_c = xgb_clf.predict_proba(X_test_clf)[:, 1]

acc_xgb_c = accuracy_score(y_test_clf, y_pred_xgb_c)
f1_xgb_c = f1_score(y_test_clf, y_pred_xgb_c)
auc_xgb_c = roc_auc_score(y_test_clf, y_prob_xgb_c)
print(f"XGBoost Classifier:   Accuracy = {acc_xgb_c:.4f}, F1 = {f1_xgb_c:.4f}, ROC-AUC = {auc_xgb_c:.4f}")
\end{codebox}


\section{Evaluation / Expected Output}
\begin{table}[htbp]
  \centering
  \caption{XGBoost classification on Heart Disease (76 test samples).}
  \begin{tabular}{|l|r|}
    \hline
\rowcolor{mlblue!8}    \textbf{Metric} & \textbf{Value} \\
    \hline
    Accuracy & 0.7763 \\
\hline
    F1-score & 0.8132 \\
\hline
    ROC-AUC & 0.8509 \\
    \hline
  \end{tabular}
\end{table}

XGBoost matches RFC on accuracy with a slightly higher F1 but a slightly lower AUC on this
small test set; differences of 0.01--0.02 are within noise for 76 samples. On the larger
regression benchmark, gradient boosting clearly leads (Experiment 8 comparison), which is the
usual pattern for tabular data.

\section{Viva / Discussion Questions}
\begin{enumerate}
  \item \textbf{Bagging vs boosting?}\par\nobreak Bagging trains models in parallel on bootstrap samples and
  combines them to reduce variance; boosting trains models sequentially so each corrects the
  previous ensemble's errors, mainly reducing bias.
  \item \textbf{What does the learning rate do?}\par\nobreak It shrinks each tree's contribution; smaller values
  need more trees but usually generalize better, larger values fit faster and can overfit.
  \item \textbf{Why can boosting overfit?}\par\nobreak Later trees keep fitting remaining errors, including label
  noise, so the ensemble can become too complex; depth limits, regularization, subsampling and
  early stopping control this.
\item \textbf{How do the learning rate and the number of trees interact?}\par\nobreak A smaller learning rate shrinks each tree's contribution and usually needs more trees; the pair must be tuned together to balance underfitting and overfitting.
\end{enumerate}


% ---------------------------------------------------------------------------
\chapter{AdaBoost}\label{exp:11}

\runninghead{Experiment 11: AdaBoost}

\quickfacts
  {Dataset}{Heart Disease (same stratified 75/25 split)}
  {Key parameters}{100 decision stumps, learning rate 0.1}
  {Headline result}{Accuracy 0.7895; F1 0.8182; ROC-AUC 0.8714 (best of the four ensembles)}

\section{Objective}
Build an adaptive boosting classifier and study how weak learners combine into a stronger
ensemble.

\section{Suggested Dataset}
Heart Disease (\texttt{data/heart\_disease.csv}), same split as the other ensembles.

\section{Required Libraries}
scikit-learn, numpy, pandas.

\section{Brief Theory}
AdaBoost trains weak learners (decision stumps) sequentially on re-weighted data.
Misclassified samples receive larger weights, so each new learner focuses on the current
errors. At round $m$ the weighted error $\epsilon_m$ determines the learner weight
\[
\alpha_m = \tfrac{1}{2}\ln\!\left(\frac{1-\epsilon_m}{\epsilon_m}\right),
\]
and the sample weights are updated as
$w_i^{(m+1)} = w_i^{(m)} \exp\!\bigl(-\alpha_m y_i h_m(x_i)\bigr)$.
The final prediction is a weighted vote of all stumps. Because the exponential loss grows
quickly for hard points, AdaBoost is sensitive to label noise and outliers.

\section{Procedure}
\begin{enumerate}
  \item Prepare the heart data and the stratified split.
  \item Fit \texttt{AdaBoostClassifier} with 100 stumps and learning rate 0.1.
  \item Predict labels and probabilities on the test set.
  \item Report accuracy, F1 and ROC-AUC; compare with RFC and XGBoost.
\end{enumerate}

\section{Reference Implementation}
% ===========================================================================
% exp11_adaboost.tex - complete code listings for Experiment 11
% Source: notebooks/04_ensemble_learning.ipynb (executed)
% ===========================================================================

\begin{codebox}[title={Core libraries}]
# ---- Core libraries ----
import numpy as np
import pandas as pd
import matplotlib.pyplot as plt
import seaborn as sns

from sklearn.datasets import fetch_california_housing
from sklearn.model_selection import train_test_split
from sklearn.ensemble import RandomForestRegressor, RandomForestClassifier, AdaBoostRegressor, AdaBoostClassifier
from sklearn.metrics import r2_score, root_mean_squared_error, mean_absolute_error, accuracy_score, f1_score, roc_auc_score

import xgboost as xgb
import catboost as cb

# Styling setup
plt.style.use('seaborn-v0_8-whitegrid' if 'seaborn-v0_8-whitegrid' in plt.style.available else 'default')
plt.rcParams['font.sans-serif'] = 'DejaVu Sans'
plt.rcParams['figure.figsize'] = (10, 6)
plt.rcParams['figure.dpi'] = 110
np.random.seed(42)
\end{codebox}

\begin{codebox}[title={Load the Heart Disease dataset (real-world classification benchmark)}]
# ---- Load the Heart Disease dataset (real-world classification benchmark) ----
df_clf = pd.read_csv('../data/heart_disease.csv')
print(f"Heart Disease dataset: {df_clf.shape[0]} rows x {df_clf.shape[1]} columns")

# Separate target and features (handle either 'target' or 'Target' as the column name)
target_col = 'target' if 'target' in df_clf.columns else 'Target'
y_clf = df_clf[target_col].values
X_clf_raw = df_clf.drop(columns=[target_col])

# One-hot encode categorical variables for RFC / XGBoost / AdaBoost
X_clf_encoded = pd.get_dummies(X_clf_raw, drop_first=True)

# Stratified split keeps the disease/no-disease ratio identical in train and test
X_train_clf, X_test_clf, y_train_clf, y_test_clf = train_test_split(
    X_clf_encoded, y_clf, test_size=0.25, random_state=42, stratify=y_clf
)

print(f"Classification Training Set: {X_train_clf.shape}, Test Set: {X_test_clf.shape}")

# ---- Random Forest Classification (majority vote of trees) ----
rfc = RandomForestClassifier(n_estimators=120, max_depth=8, random_state=42, n_jobs=-1)
rfc.fit(X_train_clf, y_train_clf)

# Predictions and probability of the positive (disease) class
y_pred_rfc = rfc.predict(X_test_clf)
y_prob_rfc = rfc.predict_proba(X_test_clf)[:, 1]

# ---- Evaluate: accuracy, F1 and ROC-AUC ----
acc_rfc = accuracy_score(y_test_clf, y_pred_rfc)
f1_rfc = f1_score(y_test_clf, y_pred_rfc)
auc_rfc = roc_auc_score(y_test_clf, y_prob_rfc)

print(f"Random Forest Classifier: Accuracy = {acc_rfc:.4f}, F1 = {f1_rfc:.4f}, ROC-AUC = {auc_rfc:.4f}")
\end{codebox}

\begin{codebox}[title={AdaBoost classification}]
# ---- AdaBoost classification ----
ada_clf = AdaBoostClassifier(n_estimators=100, learning_rate=0.1, random_state=42)
ada_clf.fit(X_train_clf, y_train_clf)

y_pred_ada_c = ada_clf.predict(X_test_clf)
y_prob_ada_c = ada_clf.predict_proba(X_test_clf)[:, 1]

acc_ada_c = accuracy_score(y_test_clf, y_pred_ada_c)
f1_ada_c = f1_score(y_test_clf, y_pred_ada_c)
auc_ada_c = roc_auc_score(y_test_clf, y_prob_ada_c)
print(f"AdaBoost Classifier:  Accuracy = {acc_ada_c:.4f}, F1 = {f1_ada_c:.4f}, ROC-AUC = {auc_ada_c:.4f}")
\end{codebox}


\section{Evaluation / Expected Output}
\begin{table}[htbp]
  \centering
  \caption{AdaBoost classification on Heart Disease (76 test samples).}
  \begin{tabular}{|l|r|}
    \hline
\rowcolor{mlblue!8}    \textbf{Metric} & \textbf{Value} \\
    \hline
    Accuracy & \best{0.7895} \\
\hline
    F1-score & \best{0.8182} \\
\hline
    ROC-AUC & \best{0.8714} \\
    \hline
  \end{tabular}
\end{table}

AdaBoost with shallow stumps slightly outperforms the deeper models on this small dataset ---
a common outcome when the signal is mostly additive and the sample size is limited. It remains
a cheap, strong baseline; its weakness is sensitivity to label noise, where regularized
gradient boosting is more robust.

\section{Viva / Discussion Questions}
\begin{enumerate}
  \item \textbf{Why use weak learners?}\par\nobreak Individually they are simple and low-variance; boosting
  combines many of them so the ensemble can fit complex boundaries while each step stays
  simple.
  \item \textbf{How does AdaBoost focus on difficult examples?}\par\nobreak After each round, weights of
  misclassified samples increase, so the next learner optimizes a distribution that emphasizes
  current errors.
  \item \textbf{What is the effect of noisy labels?}\par\nobreak Mislabeled points look ``hard'', receive
  ever-larger weights and can dominate training, degrading accuracy --- AdaBoost is known to be
  noise-sensitive.
\item \textbf{What is the difference between AdaBoost and gradient boosting?}\par\nobreak AdaBoost reweights misclassified samples and combines weak learners by weighted votes; gradient boosting fits each new learner to the residual (gradient) of the loss, which generalizes to many loss functions.
\end{enumerate}


% ---------------------------------------------------------------------------
\chapter{CatBoost}\label{exp:12}

\runninghead{Experiment 12: CatBoost}

\quickfacts
  {Dataset}{Heart Disease (raw table with categorical features: sex, cp, fbs, restecg, exang, slope, ca, thal)}
  {Key parameters}{150 iterations, depth 5, learning rate 0.08, native \texttt{cat\_features}}
  {Headline result}{Accuracy 0.7632; F1 0.7955; ROC-AUC 0.8551}

\section{Objective}
Train a gradient-boosted decision-tree model with native handling of categorical
variables.

\section{Suggested Dataset}
Heart Disease (\texttt{data/heart\_disease.csv}) using the raw (unencoded) table so CatBoost
can process categorical columns directly.

\section{Required Libraries}
catboost, pandas, scikit-learn.

\section{Brief Theory}
CatBoost is a gradient-boosting library designed for tables that mix numerical and categorical
features. It uses \textbf{ordered boosting} --- leaf values are computed on random permutations
of the data --- to avoid target leakage and prediction shift, and builds \textbf{oblivious
(symmetric) trees}, where the same split condition is used at every node of a depth level.
Categorical features are encoded internally with ordered target statistics, so the categorical
columns must be declared through \texttt{cat\_features} and the target must never leak into the
encoding.

\section{Procedure}
\begin{enumerate}
  \item Identify the categorical columns in the raw heart table and cast them to string.
  \item Use the same stratified 75/25 split on the raw features.
  \item Fit \texttt{CatBoostClassifier} with 150 iterations, depth 5, learning rate 0.08 and
        the categorical columns declared.
  \item Predict labels and probabilities; report accuracy, F1 and ROC-AUC.
  \item Compare all four classifiers on the same split.
\end{enumerate}

\section{Reference Implementation}
\input{code/exp12_catboost}

\section{Evaluation / Expected Output}
\begin{table}[htbp]
  \centering
  \caption{Ensemble classification comparison on the same Heart Disease split (76 test samples).}
  \begin{tabular}{|l|r|r|r|}
    \hline
\rowcolor{mlblue!8}    \textbf{Model} & \textbf{Accuracy} & \textbf{F1} & \textbf{ROC-AUC} \\
    \hline
    Random Forest (RFC) & 0.7763 & 0.8046 & 0.8704 \\
\hline
    XGBoost & 0.7763 & 0.8132 & 0.8509 \\
\hline
    AdaBoost & \best{0.7895} & \best{0.8182} & \best{0.8714} \\
\hline
    CatBoost & 0.7632 & 0.7955 & 0.8551 \\
    \hline
  \end{tabular}
\end{table}

CatBoost performs at a similar level to the other ensembles without any manual one-hot
encoding --- its native categorical handling does the encoding internally. On this tiny
dataset the advanced target statistics cannot show their full advantage (they matter most with
high-cardinality categories, e.g.\ thousands of city or user IDs); all differences are within
noise for 76 test samples.

\section{Viva / Discussion Questions}
\begin{enumerate}
  \item \textbf{How does CatBoost treat categorical features?}\par\nobreak It computes ordered target statistics
  (the mean target per category on random permutations, with prior smoothing) rather than
  one-hot vectors.
  \item \textbf{Why should categories not be encoded using target leakage?}\par\nobreak Plain target encoding uses
  each row's own target when encoding that row, which leaks label information and inflates
  training performance; ordered statistics avoid this.
  \item \textbf{CatBoost vs ordinary one-hot encoding?}\par\nobreak One-hot explodes dimensionality for
  high-cardinality features and loses category statistics; CatBoost encodes compactly and
  exploits target--category relationships with leakage control.
\item \textbf{Why does CatBoost often need less hyperparameter tuning?}\par\nobreak Its ordered boosting and ordered target statistics reduce target leakage, and symmetric (oblivious) trees are strong defaults, so good results are obtained with modest tuning.
\end{enumerate}

